% ECONOMICS_PROBLEM: {"id": "G12-270", "title": "Risk versus mispricing in return factors", "jel_code": "G12", "source_id": "OP-0314"}
% !TeX program = xelatex
\documentclass[11pt,a4paper]{article}
\usepackage[margin=23mm]{geometry}
\usepackage{fontspec}
\setmainfont{Latin Modern Roman}
\usepackage{amsmath,amssymb,mathtools}
\usepackage{xcolor,enumitem,booktabs,longtable,array,xurl,hyperref}
\definecolor{muted}{HTML}{53636C}
\hypersetup{colorlinks=true,urlcolor=blue,linkcolor=blue}
\setlength{\parindent}{0pt}
\setlength{\parskip}{5pt}
\newcommand{\R}{\mathbb R}
\newcommand{\E}{\mathbb E}
\newcommand{\N}{\mathbb N}
\newcommand{\ind}{\mathbf 1}
\newcommand{\Var}{\operatorname{Var}}
\newcommand{\Cov}{\operatorname{Cov}}
\newcommand{\supp}{\operatorname{supp}}
\DeclareMathOperator*{\argmin}{arg\,min}
\DeclareMathOperator*{\argmax}{arg\,max}
\newcommand{\entrymeta}[3]{{\sffamily\small\color{muted}\textbf{\texttt{#1}}\quad|\quad #2\par #3\par}}
\newcommand{\Problemref}[1]{\texttt{#1}}
\newcommand{\JELref}[2]{\texttt{#2} (JEL \texttt{#1})}
\begin{document}
\section*{Risk versus mispricing in return factors}
\textbf{Problem ID:} \texttt{G12-270}\quad \textbf{JEL:} \texttt{G12}\par
% BEGIN ECONOMICS STATEMENT
\label{problem:OP-0314}
{\sffamily\small\color{muted}Original subject group: Asset pricing and financial markets\par}
\label{definitions:problem:30007596}
\textbf{Audit classification:} Background; proposed extension. The cited paper studies factor replication. It does not verify this exact risk-versus-mispricing question as open; its replication conclusions are existing results. The cost-adjusted discrepancy is repaired.
\subsection*{Definitions and precise research target}
Fix a preregistered universe of assets, dates, predictive characteristics and executable portfolio rules. For characteristic \(k\), let \(r^k_{t+1}\) be the pre-cost excess return of its self-financing portfolio and \(c^k_{t+1}\) its trading, borrowing and market-impact cost. Let \(M_{t+1}>0\) be a discount factor generated by specified preferences and risks, with independently testable restrictions. Define the discounted net-profitability statistic
\[\alpha_k=E[M_{t+1}(r^k_{t+1}-c^k_{t+1})].\]
Even correctly priced gross portfolios can have negative \(\alpha_k\) because costs are positive; such a value is not evidence of mispricing. The research task is to determine which persistent predictive patterns reflect covariance with economically specified risks and which retain economically material discrepancies from their explicitly cost-adjusted risk-model benchmark after feasible implementation. Identify restrictions that distinguish risk-based models from belief-based mispricing models using portfolio holdings, expectations, quantities and exogenous shocks. Without such restrictions the discount factor can be chosen to fit returns, so the requested decomposition is not identified merely by a significant factor regression. Report the identified sets for the pre-cost discrepancy \(\delta_k=E[M_{t+1}r^k_{t+1}]\) and net-profitability statistic \(\alpha_k\), variation across markets and investment capacities, and correction for selection among candidate characteristics. Use genuinely subsequent or otherwise untouched observations for validation. Reproducibility of a return pattern, gross profitability and a unique economic explanation are three separate conclusions.

\textbf{Additional formal definitions and scope (editorial).} A self-financing excess payoff is the payoff of a portfolio with zero initial net cost; report the financing convention converting payoffs into returns per unit of collateral or gross capital. The pre-cost pricing restriction is \(E[M r^k]=0\) under the specified frictionless model. Consequently \(\alpha_k=-E[M c^k]\) for an otherwise correctly priced implementable portfolio; comparing \(\alpha_k\) with zero would misclassify transaction costs. Define the risk-model discrepancy as \(\delta_k=E[M r^k]\) and separately the economically attainable net-profit distribution. A positive \(\delta_k\) rejects that particular discount-factor model, while mispricing relative to all admissible rational models requires rejection throughout a prespecified class \(\mathcal M\). Restrict \(\mathcal M\) using independent consumption, holdings, or risk data: permitting every positive fitted discount factor generally does not identify a risk-versus-belief decomposition.

For the identification part, declare an admissible class \(\Theta\) of structural models, including preferences, technologies, shock laws, measurement/sampling rules and any equilibrium selector. If \(O\) denotes the complete observable history and \(T(\theta)\) the target defined above, the identified set is \(\mathcal I_T=\{T(\theta):\theta\in\Theta,\ \mathcal L_\theta(O)=\mathcal L_{\rm obs}(O)\}\); point identification means this set is a singleton. A \(do(a)\) comparison replaces stated policy/technology equations while fixing the declared exogenous law and re-solving the remaining equations. These definitions do not assert that an identification theorem holds without further restrictions.
\subsection*{Variants and assumption changes}
\begin{enumerate}
\item \textbf{Specified rational kernel (editorial).} Fix a consumption-based \(M\) independently of factors and test \(E[Mr^k]=0\) while separately estimating costs.
\item \textbf{Competing mechanisms (editorial).} Compare a prespecified rational-kernel class with a survey-disciplined belief model using untouched returns and holdings; report overlap of their observable implications.
\end{enumerate}
\subsection*{References and source audit}
\begin{enumerate}
\item Publisher verifies Jensen, Kelly and Pedersen, title, May 26, 2023 publication and DOI. Locator: Publisher abstract. Support: Factor replication results; no exact open decomposition. Abstract/metadata inspected; full body not verified. URL: \url{https://onlinelibrary.wiley.com/doi/abs/10.1111/jofi.13249}.
\end{enumerate}
\begin{enumerate}\raggedright
\item\hypertarget{definitions:source:30007596:1}{} Theis Ingerslev Jensen; Bryan Kelly; Lasse Heje Pedersen. 2023. \emph{Is There a Replication Crisis in Finance?}. Publisher abstract.
\url{https://onlinelibrary.wiley.com/doi/abs/10.1111/jofi.13249}.
DOI: \url{https://doi.org/10.1111/jofi.13249}.
\end{enumerate}

\subsection*{Common conventions: Source mathematical conventions}

These conventions apply to each entry unless explicitly replaced.

\textbf{Models and identification.} A primitive model \(m\) specifies all probability laws, feasible choices, information, equilibrium and measurement rules used by its target. The admissible class \(\mathcal M\) is fixed independently of outcomes. If \(P_O(m)\) is its observable probability law and \(\tau(m)\) the target, its identified set at observable law \(P\) is
\[
 \mathcal I_\tau(P)=\{\tau(m):m\in\mathcal M,\ P_O(m)=P\}.
\]
Point identification means that this set is a singleton. An empty set signals incompatibility of data and assumptions. Coordinate bounds need not describe the joint set. Bounds are sharp only if every retained target is attainable or arbitrarily approachable by compatible models. Finite bounds require explicit restrictions on unobserved quantities; unrestricted confounding, hidden wealth, extrapolation or arbitrary equilibrium selection can yield uninformative sets.

\textbf{Causal interventions.} A potential outcome \(Y_i(p)\) is the outcome under a complete policy regime \(p\), including timing, financing, information and market spillovers. The target population and units are fixed before treatment. With interference, use the complete assignment vector or a justified exposure mapping; individual no-interference notation alone cannot identify an economy-wide intervention. A difference of mean potential outcomes does not require the cross-world joint distribution. Conditioning on post-treatment migration, survival, enrollment or employment changes the estimand and needs separate justification.

\textbf{Statistical statements.} An estimator, confidence set, forecast or test specifies a sampling law, model class and loss/coverage criterion. Uniform \((1-\alpha)\)-coverage means \(\inf_{m\in\mathcal M}\Pr_m\{\tau(m)\in C_n\}\geq1-\alpha\), or an explicitly asymptotic version. The whole parameter space trivially has coverage, so informativeness requires a stated width, risk or power objective. A forecasting improvement is relative to a named benchmark, information set and evaluation population.

\textbf{Equilibrium and optimization.} An equilibrium comprises feasible plans, optimality, beliefs where needed, budget constraints and all market/resource clearing. The primitives or a stated benchmark define each correspondence \(\mathcal E(p;m)\). Nonemptiness is assumed or proved, never inferred just from compact policy sets. A selected-equilibrium target uses a declared selection rule; a robust target quantifies over all permitted equilibria. Use suprema or infima unless attainment is established. An \(\varepsilon\)-optimal policy is feasible and within \(\varepsilon\) of the finite optimum value. Report an empty feasible set separately.

\textbf{Welfare and accounting.} Normative utility, interpersonal normalization, social weights, numeraire, discounting and the use of public receipts are primitives. Behavioral utility need not coincide with the welfare criterion. Household welfare includes ownership income and rents once; adding profits again to compensating variation generally double-counts them. A monetary equivalent is defined by an explicit transfer and price convention, with monotonicity and range conditions. Average effects, mechanism contributions, transfers and welfare losses are different targets. Infinite sums require convergence and dynamic feasible plans require terminal or no-Ponzi conditions.

\textbf{Source versions.} A working-paper date, revision date and journal publication year are distinguished. A DOI for a journal version is not automatically the DOI of an earlier manuscript. Locators refer to the stated version and printed pages where specified. A metadata-only check does not verify a claim about a full-text conclusion. Every entry's source subsection narrows the provenance claim accordingly.
% END ECONOMICS STATEMENT
\end{document}
